\documentclass[../../main]{subfiles}
\zexternaldocument*{../1.Category-Theory/exported/aux/main}
\begin{document}
\chapter{Old}
\section{Complexes and (co)homology}

From now on, we will be working in a fixed \nameref{def:abelian-category} \(\mathsf{A}\).
By the \nameref{thm:freyd-mitchell} we can work with objects and morphisms using set-theoretic tools when needed.
\begin{definition}[Image and coimage] \label{def:image}
  Let \(\varphi: A \to B\) be a morphism in an \nameref{def:abelian-category}.

  The \textbf{image} of \(\varphi\) is \(\im \varphi = \ker(\coker \varphi)\) and the \textbf{coimage} of \(\varphi\) is \(\coim \varphi = \ker(\coker \varphi)\).
\end{definition}


\begin{definition}[Cohomology] \label{def:cohomology}
  Let \(C^{\bullet}\) be a \nameref{def:complex} in an \nameref{def:abelian-category} \(\mathsf{A}\).
  We define the \textbf{\(i\)-th cohomology object} of \(C^{\bullet}\) as
  \begin{equation*}
    \operatorname{H}^{i}(C^{\bullet}) = \quot{\ker d^{i}}{\im d^{i-1}}
  \end{equation*}
\end{definition}

\begin{definition}[Category of complexes] \label{def:complex-category}
  Let \(\mathsf{A}\) be an \nameref{def:abelian-category}.
  The \textbf{\nameref{def:category} of \namerefs{def:complex}} \(\mathsf{C}(\mathsf{A})\) consists of:
  \begin{enumerate}
    \item \hyperref[def:complex]{Cochain complexes} in \(\mathsf{A}\) as objects.
    \item For \((M^{\bullet}, d_{M^{\bullet}}^{\bullet}), (N^{\bullet}, d_{N^{\bullet}}^{\bullet}) \in \mathsf{C}(\mathsf{A})\), \(\Hom_{\mathsf{C}(\mathsf{A})}(M^{\bullet}, N^{\bullet})\) consists of the commutative diagrams of the form
          \begin{center}
            \begin{tikzcd}
              \cdots \arrow[r] & M^{i-1} \arrow[r, "d_{M^{\bullet}}^{i-1}"] \arrow[d, "\alpha^{i-1}"'] & M^{i} \arrow[r, "d_{M^{\bullet}}^{i}"] \arrow[d, "\alpha^{i}"'] & M^{i+1} \arrow[r, "d_{M^{\bullet}}^{i+1}"] \arrow[d, "\alpha^{i+1}"'] & \cdots \\
              \cdots \arrow[r] & N^{i-1} \arrow[r, "d_{N^{\bullet}}^{i-1}"] & N^{i} \arrow[r, "d_{N^{\bullet}}^{i}"] & N^{i+1} \arrow[r, "d_{N^{\bullet}}^{i+1}"] & \cdots
            \end{tikzcd}
          \end{center}
          We denote by \(\alpha^{\bullet}: M^{\bullet} \to N^{\bullet}\) the morphism determined by the collection of \(\alpha^{i}\).
  \end{enumerate}
\end{definition}

\begin{lemma}
  The \nameref{def:complex-category} of an \nameref{def:abelian-category} is an \nameref{def:abelian-category}.
\end{lemma}

\begin{lemma}
  Let \(\mathsf{A}\) be a \nameref{def:abelian-category} and let \(\mathsf{C}(\mathsf{A})\) be its \nameref{def:complex-category}.

  Then for every integer \(i\), the assignement \(\operatorname{H}^{i}: M^{\bullet} \mapsto H^{i}(M^{\bullet})\) taking a \nameref{def:complex} \(M^{\bullet}\) to its \(i\)-th \nameref{def:cohomology} is a \nameref{def:functor} \(\mathsf{C}(\mathsf{A}) \to \mathsf{A}\).
\end{lemma}
\begin{proof}
  We want to show that for all \namerefs{def:complex} \(M^{\bullet}, N^{\bullet} \in \mathsf{C}(\mathsf{A})\) and \(i \in \mathbb{Z}\), \(\operatorname{H}^{i}\) induces in a natural way a homomorphism of abelian groups
  \begin{equation*}
    \Hom_{\mathsf{C}(\mathsf{A})}(M^{\bullet}, N^{\bullet}) \to \Hom_{\mathsf{A}}(\operatorname{H}^{i}(M^{\bullet}), \operatorname{H}^{i}(N^{\bullet}))
  \end{equation*}

  Let \(\alpha^{\bullet}: M^{\bullet} \to N^{\bullet}\) be a morphism in \(\mathsf{C}(\mathsf{A})\). Let's consider the relevant of the commutative diagram defining \(\alpha^{\bullet}\):
  \begin{center}
    \begin{tikzcd}
      M^{i-1} \arrow[r, "d_{M^{\bullet}}^{i-1}"] \arrow[d, "\alpha^{i-1}"'] & M^{i} \arrow[r, "d_{M^{\bullet}}^{i}"] \arrow[d, "\alpha^{i}"'] & M^{i+1} \arrow[d, "\alpha^{i+1}"'] \\
      N^{i-1} \arrow[r, "d_{N^{\bullet}}^{i-1}"] & N^{i} \arrow[r, "d_{N^{\bullet}}^{i}"] & N^{i+1}
    \end{tikzcd}
  \end{center}
  By the commutativity of the right square we have that \(d_{N^{\bullet}}^{i} \circ \alpha^{i} = \alpha^{i+1} \circ d_{M^{\bullet}}^{i}\) is 0 on \(\ker d^{i}\).
  Therefore the restriction of \(\alpha^{i}\) to \(\ker d_{M^{\bullet}}^{i}\) factors through \(\ker d_{N^{\bullet}}^{i}\).
  Composing with the projection we get a morphism
  \begin{equation*}
    \ker d_{M^{\bullet}}^{i} \to \operatorname{H}^{i}(N^{\bullet})
  \end{equation*}
  Since this morphism is 0 when restricted to \(\im d_{M^{\bullet}}^{i-1}\) by the commutativity of the left square we get that \(\alpha^{i}\) induces a morphism
  \begin{equation*}
    \operatorname{H}^{i}(M^{\bullet}) \to \operatorname{H}^{i}(N^{\bullet})
  \end{equation*}
  in \(\mathsf{A}\) as we wanted.
\end{proof}

\begin{remark}
  We can view \nameref{def:cohomology} as a \nameref{def:functor} \(\mathsf{C}(\mathsf{A}) \to \mathsf{C}(\mathsf{A})\) by placing each \nameref{def:cohomology} object \(\operatorname{H}^{i}(M^{\bullet})\) in degree \(i\) and connecting them using 0 morphisms, obtaining a complex \(\operatorname{H}^{\bullet}(M^{\bullet})\).

  Note that the \nameref{def:cohomology} of this particular \nameref{def:complex} verifies
  \begin{equation*}
    \operatorname{H}^{\bullet}(\operatorname{H}^{\bullet}(M^{\bullet})) = \operatorname{H}^{\bullet}(M^{\bullet})
  \end{equation*}
\end{remark}




\end{document}
